\lecture{19}
\title{Non-parametric Tests}

\begin{document}

So far, we've only looked at tests for parameters of a distribution.
Non-parametric tests refer to tests about the distribution as a whole;
the three examples of such tests we'll discuss in this lecture are:
\begin{itemize}
    \item \textbf{Kolmogorov-Smirnov.}
    Given $X_1, \dots, X_n$ sampled (i.i.d.)
    from a distribution with CDF $F$, we have:
    \[ H_0 := \text{``The CDF $F$ matches another CDF $F_0$.''}
    ~~~ \parallel ~~~
    H_1 := \text{``The CDF $F$ does not match another CDF $F_0$.''} \]
    \item \textbf{Kolmogorov-Lilliefors.}
    Given $X_1, \dots, X_n$ sampled (i.i.d.)
    from a distribution $X \sim \mathbb{P}$, we have:
    \[ H_0 := \text{``The random variable $X$ is Gaussian.''}
    ~~~ \parallel ~~~
    H_1 := \text{``The random variable $X$ is not Gaussian.''} \]
    \item \textbf{Two-Sample Test.}
    Given $X_1, \dots, X_n \sim F_X$
    and $Y_1, \dots, Y_m \sim F_Y$ (i.i.d.),
    where $F_X$ and $F_Y$ are CDFs, we have:
    \[ H_0 := \text{``The CDF $F_X$ matches the CDF $F_Y$.''}
    ~~~ \parallel ~~~
    H_1 := \text{``The CDF $F_X$ does not match the CDF $F_Y$.''} \]
\end{itemize}
All three of these tests share the same idea behind their test statistic.

\textbf{Definition.}
Given $X_1, \dots, X_n$ sampled (i.i.d.) from a distribution,
their \textbf{empirical CDF} is
$\hat{F}_n(t) := \frac{1}{n} \sum_{i = 1}^n \mathds{1}(X_i \leq t)$.

\begin{center}
    \frame{\includegraphics[width=0.4\textwidth]{lecture-19/empirical-cdf}}
\end{center}

Then the \textbf{test statistic} for all three tests will look something like
$T_n := \sup_t \left|\hat{F}_n(t) - F_0(t)\right|$,
which represents the greatest absolute difference
between the empirical CDF $\hat{F}_n(t)$ and the expected CDF $F_0(t)$.

\subsection{The Kolmogorov-Smirnov Test}

Here's the main idea behind the test.

\begin{theorem}{Kolmogorov-Smirnov Distribution}
    Assuming the null hypothesis,
    suppose $X_1, \dots, X_n$ are sampled (i.i.d.)
    from the distribution with continuous CDF $F_0$.
    Then the distribution of the test statistic $T_n$
    does \emph{not} depend on $F_0$.
\end{theorem}

\begin{remark}
    This should not be particularly surprising.
    For example, it should be clear why
    $F_0(t)$ and $F_0(\frac{t}{2})$ would yield the same test statistic:
    scaling the $x$-axis of a CDF shouldn't affect
    the maximum-detected difference in $y$-values.
\end{remark}

\begin{proof}
    The point is that
    the CDF $F_0$ has no relevance at all,
    and the lemma below explains why.
    \begin{quote}
        \textbf{Lemma.}
        If a random variable $X$ has continuous CDF $F_0$,
        then $F_0(X) \sim \mathrm{Unif}(0, 1)$.

        \emph{Proof:}
        Ignoring the real-analysis technicalities,
        here's the gist of why this should feel true.

        When we say $X$ has continuous CDF $F_0$,
        that just means $X$ can be generated by
        picking a random number $\lambda \sim \mathrm{Unif}(0, 1)$,
        then setting $X = F_0^{-1}(\lambda)$.
        The conclusion is immediate from here. ~ $\square$
    \end{quote}
    The above lemma lets us rewrite the empirical distribution,
    which then lets us write the test statistic $T_n$
    with no reference to $F_0$.
    Below, we set $U_i := F_0(X_i)$, so that by the lemma,
    these are all i.i.d. samples from $\mathrm{Unif}(0, 1)$.
    \begin{align*}
        T_n = \sup_{t \in \mathbb{R}} \left|\hat{F}_n(t) - F_0(t)\right|
        & = \sup_{t \in \mathbb{R}} \left|\left[
            \frac{1}{n}\sum_{i = 1}^n \mathds{1}(X_i \leq t)
        \right] - F_0(t)\right| \\
        & = \sup_{t \in \mathbb{R}} \left|\left[
            \frac{1}{n}\sum_{i = 1}^n \mathds{1}(U_i \leq F_0(t))
        \right] - F_0(t)\right|
        = \sup_{s \in [0, 1]} \left|\left[
            \frac{1}{n}\sum_{i = 1}^n \mathds{1}(U_i \leq s)
        \right] - s\right| \sim \mathrm{KS}_n.
    \end{align*}
    That last expression is the definition of the
    \textbf{Kolmogorov-Smirnov distribution} $\mathrm{KS}_n$,
    which does not depend on $F_0$. ~ $\blacksquare$
\end{proof}

And, of course, the above theorem
easily gives rise to a Kolmogorov-Smirnov Test
of level $\alpha$.
Just compute the test statistic $T_n$
and reject $H_0$ if $T_n > \mathrm{KS}_{n, \alpha}$.

\subsection{The Kolmogorov-Lilliefors Test}

To test whether a random variable is Gaussian,
we rework our test statistic like so:
\[
    T_n := \sup_t\left|
        \hat{F}_n(t) - \Phi_{\hat{\mu}, \hat{\sigma}^2}(t)
    \right| ~~~ \text{ where } ~~~
    \hat{\mu} = \text{sample mean} ~~~ \text{ and } ~~~
    \hat{\sigma}^2 = \text{sample variance}
\]

\begin{theorem}{Kolmogorov-Lilliefors Distribution}
    Assuming the null hypothesis,
    suppose $X_1, \dots, X_n$ are sampled (i.i.d.)
    from the distribution $\mathcal{N}(\mu, \sigma^2)$.
    Then the distribution of the test statistic $T_n$
    does not depend on $\mu$ or $\sigma^2$.
\end{theorem}

\begin{remark}
    This one \emph{really} shouldn't be unexpected.
    It should be clear why the CDFs $F_0(t)$ and
    $F_0\left(\frac{t - \mu}{\sigma}\right)$
    would yield the same test statistic,
    because scaling or shifting the $x$-axis
    shouldn't affect the maximum-detected
    difference in $y$-values.
\end{remark}

\begin{proof}
    Say that $X_i = \mu + \sigma \cdot Z_i$,
    where $Z_1, \dots, Z_n \sim \mathcal{N}(0, 1)$.
    The goal is to rewrite $T_n$
    so that its value is dependent only on
    the values of $\{Z_i\}_{i = 1}^n$
    (which, of course, do not depend on $\mu$ or $\sigma^2$).

    Here are the relevant lemmas.
    Their proofs are very straightforward.

    \begin{quote}
        \textbf{Lemma 1.}
        Let $\hat{G}_n(t)$ denote the empirical CDF of the $\{Z_i\}_{i = 1}^n$.
        Then $\hat{F}_n(t) = \hat{G}_n\left(
            \left[\frac{t - \mu}{\sigma}\right]
        \right)$.

        \textbf{Lemma 2.}
        We have $\hat{\sigma}^2
        = \sigma^2 \cdot \hat{\sigma}_{Z, n}^2$,
        where $\hat{\sigma}_{Z, n}$ is
        a function of only the $\{Z_i\}_{i = 1}^n$.

        \textbf{Lemma 3.}
        As a corollary of Lemma 2,
        if $\Phi$ is the CDF of $\mathcal{N}(0, 1)$,
        then we have:
        \[ \Phi_{\hat{\mu}, \hat{\sigma}^2}(t)
        = \Phi\left(\frac{t - \bar{X}_n}{\hat{\sigma}}\right)
        = \Phi\left(\frac{
                t - \mu - \sigma \cdot \bar{Z}_n
            }{
                \sigma \cdot \hat{\sigma}_{Z, n}
            }\right)
        = \Phi\left(
            \frac{\left[\frac{t - \mu}{\sigma}\right]
            - \bar{Z}_n}{\hat{\sigma}_{Z, n}}
        \right). \]
    \end{quote}
    Using Lemma 1 and Lemma 3,
    we can now rewrite $T_n$ like so:
    \[
        T_n = \sup_{t \in \mathbb{R}} \left|
            \hat{F}_n(t) - \Phi_{\hat{\mu}, \hat{\sigma}^2}(t)
        \right|
        = \sup_{t \in \mathbb{R}} \left|
            \hat{G}_n\left(
                \left[\frac{t - \mu}{\sigma}\right]
            \right) -
            \Phi\left(
                \frac{\left[\frac{t - \mu}{\sigma}\right]
                - \bar{Z}_n}{\hat{\sigma}_{Z, n}}
            \right)
        \right|
        = \sup_{s \in \mathbb{R}} \left|
            \hat{G}_n(s) -
            \Phi\left(\frac{s - \bar{Z}_n}{\hat{\sigma}_{Z, n}}\right)
        \right| \sim \mathrm{KL}_n.
    \]
    The RHS of the above is dependent only on the $\{Z_i\}_{i = 1}^n$,
    so its distribution is dependent only on $n$.
    It follows the \textbf{Kolmogorov-Lilliefors distribution} $\mathrm{KL}_n$.
    ~ $\blacksquare$
\end{proof}

And, of course, you get a Kolmogorov-Lilliefors test out of this, too.

\begin{remark}
    It turns out that $\mathrm{KL}_{n, \alpha} < \mathrm{KS}_{n, \alpha}$
    for all $n$.
    This means that if we mistakenly used the $\mathrm{KS}$ quantiles
    for a KL test, we would reject the null too infrequently.
\end{remark}

\subsection{The Two-Sample Test}

To check if two sets of samples $\{X_i\}_{i = 1}^n$ and $\{Y_i\}_{i = 1}^m$
came from the same CDF, we use the following test statistic:
\[ T_{n, m} := \sup_t \left| \hat{F}_{X, n}(t) - \hat{F}_{Y, m}(t) \right|. \]

\begin{theorem}{Two-Sample Distribution}
    Assuming the null hypothesis,
    suppose $\{X_i\}_{i = 1}^n$ and $\{Y_i\}_{i = 1}^m$
    are both sampled from the same continuous CDF $F$.
    Then the distribution of the test statistic $T_{n, m}$
    does not depend on $F$.
\end{theorem}

\begin{proof}
    Same proof as with the $\mathrm{KS}$ distribution.
    Say $U_i := F(X_i)$ and $V_i := F(Y_i)$, which by the lemma
    are i.i.d. samples from $\mathrm{Unif}(0, 1)$.  Then:
    \begin{align*}
        T_{n, m} & = \sup_{t \in \mathbb{R}}\left|
            \left[\frac{1}{n}\sum_{i = 1}^n
            \mathds{1}(X_i \leq t)\right] -
            \left[\frac{1}{m}\sum_{i = 1}^m
            \mathds{1}(Y_i \leq t)\right]
        \right| \\
        & = \sup_{t \in \mathbb{R}}\left|
            \left[\frac{1}{n}\sum_{i = 1}^n
            \mathds{1}(U_i \leq F(t))\right] -
            \left[\frac{1}{m}\sum_{i = 1}^m
            \mathds{1}(V_i \leq F(t))\right]
        \right| \\
        & = \sup_{s \in [0, 1]}\left|
            \left[\frac{1}{n}\sum_{i = 1}^n
            \mathds{1}(U_i \leq s)\right] -
            \left[\frac{1}{m}\sum_{i = 1}^m
            \mathds{1}(V_i \leq s)\right]
        \right|.
    \end{align*}
    And that doesn't depend on $F$. ~ $\blacksquare$
\end{proof}

And, yeah, you can perform tests with this, too.

\end{document}